\documentclass[11pt]{article}

\usepackage[margin=1in]{geometry}
\usepackage{amsmath,amssymb,bm}
\usepackage{microtype}
\usepackage[hidelinks]{hyperref}
\usepackage{enumitem}
\usepackage{booktabs}

\newcommand{\bk}{\mathbf{k}}
\newcommand{\bq}{\mathbf{q}}
\newcommand{\br}{\mathbf{r}}
\newcommand{\kF}{k_{\mathrm F}}
\newcommand{\qTF}{q_{\mathrm{TF}}}

\title{\vspace{-1.2cm}Graphene Minimal Conductivity from Three-Dimensional Charged-Impurity Disorder}
\author{}
\date{}

\title{Challenge 42: Final answer}
\author{}
\begin{document}
\maketitle
Therefore in summary, we have:

\begin{equation}
\boxed{
\begin{gathered}
\alpha=-\frac13,\qquad \beta=\frac23,\\[2pt]
\xi\propto n_i^{-1/3},\qquad
\Delta V_g\propto n_i^{2/3}.
\end{gathered}}
\end{equation}
For an insulating 3D TI bulk, charged impurities can create electron/hole puddles on the 2D Dirac surface and an analogous charge-neutrality plateau. Their screened Coulomb potential is long-ranged and dominated by small-$q$ scattering; for comparable disorder strength this generally gives a longer transport mean free path than short-range disorder.
\end{document}
